\documentclass[11pt]{article}

\usepackage[a4paper,margin=2.5cm]{geometry}
\usepackage{graphicx}
\usepackage{hyperref}
\usepackage{xcolor}
\usepackage{listings}
\usepackage{graphicx}
\usepackage{float}
\usepackage{subcaption}
\usepackage{tabularx}
\usepackage{parskip}
\usepackage{cleveref} 
\usepackage{booktabs}
\usepackage{titling}
\setlength{\droptitle}{-3cm}
\lstset{
    language=bash,
    basicstyle=\ttfamily\small,
    keywordstyle=\color{blue},
    commentstyle=\color{gray},
    stringstyle=\color{green!50!black},
    backgroundcolor=\color{gray!5},
    frame=single,
    rulecolor=\color{gray!40},
    breaklines=true,
    captionpos=b,
    showstringspaces=false
}

\title{\textbf{HPC - Get to know your GPU}}
\author{
Do Minh Quang - 2540095\\
University of Science and Technology of Hanoi
}
\date{}

\begin{document}

\maketitle

\section{Implementation}
The experiment was run in the GPU of Google Colab notebook. 

\begin{lstlisting}[language=Python, label={lst:code}]
import numba
from numba import cuda

cuda.detect()

# Device name
device = cuda.select_device(0)

print(f"Device id: {device.id}")
print(f"Device name: {device.name}")

# Core info
cores_per_SM_dict = {
    (2,0) : 32, (2,1) : 48, (3,0) : 192,
    (3,5) : 192, (3,7) : 192, (5,0) : 128,
    (5,2) : 128, (6,0) : 64, (6,1) : 128,
    (7,0) : 64, (7,5) : 64, (8,0) : 64,
    (8,6) : 128, (8,9) : 128, (9,0) : 128
    }

multiprocessor_count = device.MULTIPROCESSOR_COUNT
cores_per_SM = cores_per_SM_dict.get(device.compute_capability, 128)
total_cores = multiprocessor_count * cores_per_SM

print(f"Multiprocessor count: {multiprocessor_count}")
print(f"Core count: {cores_per_SM}")
print(f"Total cores: {total_cores}")

# Memory info
mem_info = cuda.current_context().get_memory_info()
print(f"Total memory: {mem_info[1]}")
\end{lstlisting}

\section{Results}

The GPU information obtained using \texttt{numba.cuda} is
summarized in \Cref{tab:gpu}.

\begin{table}[H]
\centering
\caption{GPU information reported by \texttt{numba.cuda} on Google Colab.}
\label{tab:gpu}
\begin{tabular}{lr}
\toprule
\textbf{Property} & \textbf{Value} \\
\midrule
Device ID & 0 \\
Device name & Tesla T4 \\
Compute capability & 7.5 \\
Multiprocessor count & 40 \\
Cores per SM & 64 \\
Total CUDA cores & 2,560 \\
Total memory & 15,637,086,208 bytes \\
 & $\approx$ 14.56 GiB \\
\bottomrule
\end{tabular}
\end{table}


\end{document}